% Revised 22 October 2026: final student as deployed model; source-traced macros.
\chapter*{Abstract}
\addcontentsline{toc}{chapter}{Abstract}

Multi-person human mesh recovery estimates articulated body surfaces and their positions from one image.
This thesis studies a compact Anny-based reconstruction system using a ViT-B student distilled from a larger
ViT-L teacher. Both follow the single-shot Multi-HMR framework. Anny expresses body shape through named
parameters, including age, height and proportions, and provides a common representation for the student,
teacher and evaluation.

The training pipeline combines complete synthetic supervision from AnnyOne with real images from 3DPW.
A fitting procedure converts 3DPW's SMPL annotations to Anny using geometric correspondence and virtual
torso joints. The corrected labels contain \DataThreeDPWTotalImages{} images and
\DataThreeDPWTotalPeople{} labelled person instances. Teacher training uses consistent head-joint anchors,
a fixed bone-rotation correction and a mixture containing \DataMixThreeDPW\,\% real images. A learned projector
connects the student's features to prediction heads initialized from the teacher. Ground-truth supervision
is combined with feature and output distillation losses.

The student has \StudentParams\,million parameters and requires \StudentGFLOPs\,GFLOPs in the recorded
one-person configuration, compared with \FinalParams\,million parameters and \FinalGFLOPs\,GFLOPs for the
teacher. Paired L40S forward latency is \StudentMsLFortyS\,ms for the student and \TeacherMsLFortyS\,ms for
the teacher. The reduction has an accuracy cost: student PVE is \StudentPVE\,mm on the synthetic holdout
and \StudentThreeDPW\,mm on converted 3DPW, compared with \FinalPVE\,mm and \FinalThreeDPW\,mm for the
teacher. The student's 3DPW PVE is lower than the released Multi-HMR-Anny baseline's \NAVERThreeDPW\,mm,
although training exposure differs. The baseline's synthetic-holdout result is potentially in-sample.

A frame-wise application compares the student and baseline on the same walking video. The student provides
a tracked prediction in all \VideoFrames{} frames; the baseline provides one in \VideoBaselineMatched{} frames.
These counts establish execution and tracking coverage on one clip. Pose jitter and shape consistency still
require inspection of the video and measurements. Remaining limitations include the student's accuracy gap,
adult-only synthetic holdout coverage, uncertain real-image shape and residual depth errors. All reported
3DPW errors use converted Anny labels and are not directly comparable with published SMPL-space results.
